\chapter{Dépannage}
\minitoc

\section{Problèmes au démarrage}

\subsection{« No module named cirtdefense »}

Le paquet n'est pas installé. Le code se trouve dans \texttt{src/}, et Python
ne l'y cherche pas tout seul.

\begin{lstlisting}[language=bash]
$ pip install -e ".[dev]"
\end{lstlisting}

\subsection{« WinError 10013 » sous Windows}

Le système refuse à Python le droit d'écouter sur ce port. Ce n'est pas un
problème de la plateforme. Deux causes, par ordre de fréquence.

\textbf{Plages réservées par Hyper-V, WSL ou Docker Desktop.} Ces composants
réservent dynamiquement des milliers de ports. Dans un PowerShell
administrateur :

\begin{lstlisting}[language=bash]
PS> netsh interface ipv4 show excludedportrange protocol=tcp
PS> net stop winnat
PS> cirtd serve --host 127.0.0.1 --port 8000
PS> net start winnat
\end{lstlisting}

\textbf{Antivirus ou pare-feu tiers.} Pour trancher, testez un serveur
minimal : si celui-ci échoue aussi, la plateforme est hors de cause.

\begin{lstlisting}[language=bash]
PS> python -m http.server 8000 --bind 127.0.0.1
\end{lstlisting}

\subsection{La console ne s'affiche pas, l'API répond}

Vous avez installé sans le tiret \texttt{-e}. L'interface, qui vit dans
\texttt{web/}, n'a pas été installée. Réinstallez en mode éditable.

Pour confirmer, ouvrez \texttt{/health} : le champ \texttt{web\_root\_exists}
doit valoir \texttt{true}.

\section{Problèmes de fonctionnement}

\subsection{La plateforme n'agit sur rien}

Vérifiez dans l'ordre :

\begin{enumerate}[leftmargin=*]
  \item le bandeau — l'autonomie est-elle active~?
  \item le coupe-circuit — s'est-il déclenché~?
  \item les décisions, dans le journal — quelle issue portent-elles~?
\end{enumerate}

L'issue « aucune réponse engagée : la menace n'est pas documentée » est un
\textbf{refus}, pas une panne. Elle signale que le fonds documentaire ne couvre
pas ce que vous lui envoyez.

\subsection{Les boutons de décision ne répondent pas}

Ce défaut a existé : les boutons demandaient le motif par une boîte de dialogue
du navigateur, et lorsque celui-ci les supprimait, rien ne se passait — sans le
moindre message. Il est corrigé : le motif se saisit désormais dans la page.

Si vous rencontrez encore ce comportement, vous exécutez une version
antérieure : mettez à jour.

\subsection{« database is locked »}

Deux processus écrivent en même temps dans la base. N'exécutez pas la boucle de
contrôle dans un processus séparé de l'API tant que le délai d'attente sur
verrou n'est pas configuré.

\subsection{La chaîne du journal est rompue}

Ce n'est pas une anomalie de fonctionnement mais un incident de sécurité :
une entrée a été modifiée en dehors de l'application. Préservez le fichier de
base en l'état, ne le réparez pas, et ouvrez une enquête.

\section{Limitations connues}

Elles sont énoncées ici plutôt que découvertes en exploitation.

\begin{table}[h]
\centering
\begin{tabular}{@{}p{5.6cm}p{8.4cm}@{}}
\toprule
\textbf{Limitation} & \textbf{Conséquence pratique} \\
\midrule
La boucle de contrôle ne se déclenche pas dans la topologie Docker documentée &
Elle tourne et ne vérifie rien, silencieusement. Exécutez-la dans le même
processus que l'API. \\[8pt]
L'apprentissage porte sur la qualification, pas sur la réponse &
Un type adopté est reconnu à l'occurrence suivante, mais traité par repli
faute de playbook écrit pour lui. \\[8pt]
Trois sources produisent trois clés de corrélation pour une même machine &
Un serveur vu par trois capteurs ouvre trois interventions au lieu d'une. \\[8pt]
La criticité d'un actif vaut 3 par défaut &
Sans inventaire du parc, la promotion de gravité et le score de risque sont
aplatis. \\[8pt]
Pas de délai d'attente sur verrou SQLite &
Deux écrivains simultanés produisent une erreur immédiate. \\[8pt]
Le classement des substitutions dépend du catalogue de réversibilité &
Une entrée mal déclarée produit une proposition mal classée. C'est le catalogue
qu'il faut corriger, pas le code. \\
\bottomrule
\end{tabular}
\caption{Limitations connues à la révision 1.0 de ce guide}
\end{table}

\section{Où chercher davantage}

\begin{description}[leftmargin=3.4cm, style=nextline]
  \item[\texttt{/docs}] La documentation interactive de l'API, engendrée par
        FastAPI. Toutes les routes y sont décrites et essayables.
  \item[\texttt{/health}] L'état du système en une requête : chemin du paquet,
        racine de l'interface, présence des fichiers.
  \item[\texttt{docs/}] Dans le dépôt : architecture, structure, traçabilité
        des exigences, et les décisions d'architecture (ADR).
  \item[Journal d'audit] La source de vérité sur ce que la plateforme a fait.
        Quand une explication manque, elle s'y trouve.
\end{description}
